\chapter{Bắt đầu lại}
\markboth{Bắt đầu lại}{Start Again}

\section{Tỉnh ra muộn còn hơn ngủ tiếp.}
\begin{truyen}{Tỉnh ra muộn còn hơn ngủ tiếp.}{Wake up.}
\chuhoa{C}{âu này khiến nhiều em ngẩng lên. Thầy không cần học sinh hoàn hảo ngay; điều thầy cần là các em đừng dùng chuyện “muộn rồi” để biện minh cho việc tiếp tục buông xuôi.}
\textit{(The sentence made many students look up. The teacher did not need perfection immediately; he needed students not to use “it's too late” as an excuse to keep giving up.)}
\end{truyen}

\begin{baihoc}
Nhận ra muộn vẫn còn giá trị, miễn là em chịu đổi từ bây giờ.
\textit{(Realizing late still matters, as long as you change from now on.)}
\end{baihoc}

\begin{ghinhoanh}
\textbf{Short English to remember:}

Late awakening is better than sleeping on.

\end{ghinhoanh}

\begin{tuhocnhanh}
1. Em có đang nghĩ mình “muộn quá rồi” không? \textit{Do you think you are “already too late”?}

2. Nếu chưa ngủ tiếp, em sẽ làm gì khác? \textit{If you choose not to sleep on, what will you do differently?}
\end{tuhocnhanh}
\ngancach

\section{Em không cần giỏi ngay, em cần dừng lười trước đã.}
\begin{truyen}{Em không cần giỏi ngay, em cần dừng lười trước đã.}{Wake up.}
\chuhoa{T}{hầy nói rất thật: nhiều em chưa tiến không phải vì thiếu thông minh, mà vì chưa chịu bỏ thói lười. Khi ngưng kéo lùi chính mình, việc học mới bắt đầu tiến lên được.}
\textit{(The teacher spoke plainly: many students are not stuck because they lack intelligence, but because they have not abandoned laziness. When you stop pulling yourself backward, learning can finally move forward.)}
\end{truyen}

\begin{baihoc}
Điều đầu tiên không phải là giỏi nhanh, mà là đừng tự cản mình nữa.
\textit{(The first goal is not to become good quickly, but to stop blocking yourself.)}
\end{baihoc}

\begin{ghinhoanh}
\textbf{Short English to remember:}

You don't need instant brilliance, just less self-sabotage.

\end{ghinhoanh}

\begin{tuhocnhanh}
1. Thói lười nào đang cản em nhiều nhất? \textit{Which lazy habit is holding you back the most?}

2. Em sẽ dừng nó từ đâu? \textit{Where will you start stopping it?}
\end{tuhocnhanh}
\ngancach

\section{Hối hận không đổi điểm, hành động mới đổi.}
\begin{truyen}{Hối hận không đổi điểm, hành động mới đổi.}{Wake up.}
\chuhoa{C}{ả lớp im sau câu này. Buồn, tiếc, hay tự trách có thể cho thấy em còn biết nghĩ, nhưng nếu nó không kéo theo việc làm mới thì điểm số và tương lai vẫn y nguyên.}
\textit{(The class fell silent after this line. Feeling sad, regretful, or self-critical may show that you still care, but if none of it leads to new action, grades and the future remain the same.)}
\end{truyen}

\begin{baihoc}
Cảm xúc chỉ có ích khi nó được chuyển thành việc làm cụ thể.
\textit{(Emotion becomes useful only when it is turned into concrete action.)}
\end{baihoc}

\begin{ghinhoanh}
\textbf{Short English to remember:}

Regret changes nothing; action does.

\end{ghinhoanh}

\begin{tuhocnhanh}
1. Em đang tiếc điều gì? \textit{What are you regretting?}

2. Hành động nào sẽ biến tiếc nuối thành thay đổi? \textit{What action can turn regret into change?}
\end{tuhocnhanh}
\ngancach

\section{Bắt đầu lại hôm nay vẫn sớm hơn than thở vào năm sau.}
\begin{truyen}{Bắt đầu lại hôm nay vẫn sớm hơn than thở vào năm sau.}{Wake up.}
\chuhoa{C}{âu này không làm em xấu hổ, nó mở cho em một lối thoát. Chừng nào còn bắt đầu được, chừng đó vẫn còn hy vọng thật chứ không chỉ là tự an ủi.}
\textit{(This sentence is not meant to shame you; it opens an exit. As long as you can still begin, there is still real hope rather than mere self-comfort.)}
\end{truyen}

\begin{baihoc}
Bắt đầu lại sớm luôn rẻ hơn việc kéo dài trì trệ thêm một năm.
\textit{(Starting again early is always cheaper than dragging stagnation into another year.)}
\end{baihoc}

\begin{ghinhoanh}
\textbf{Short English to remember:}

Start again now, not next year.

\end{ghinhoanh}

\begin{tuhocnhanh}
1. Nếu bắt đầu lại hôm nay, em sẽ bắt đầu từ môn hay thói quen nào? \textit{If you restart today, which subject or habit will you begin with?}

2. Điều gì xảy ra nếu em chờ thêm? \textit{What happens if you wait longer?}
\end{tuhocnhanh}
\ngancach

\section{Một giờ học thật còn hơn cả tuần giả vờ cố gắng.}
\begin{truyen}{Một giờ học thật còn hơn cả tuần giả vờ cố gắng.}{Wake up.}
\chuhoa{T}{hầy muốn lớp hiểu rằng sự thật thà với nỗ lực quan trọng hơn vẻ ngoài bận rộn. Một giờ tập trung thật có thể thay đổi nhiều hơn cả tuần cứ ngồi đó cho có cảm giác mình đang học.}
\textit{(The teacher wanted the class to understand that honesty in effort matters more than the appearance of being busy. One truly focused hour can change more than a whole week of merely sitting there to feel productive.)}
\end{truyen}

\begin{baihoc}
Chất lượng của một khoảng học đôi khi quan trọng hơn số lượng em khoe ra.
\textit{(The quality of a study session is sometimes more important than the amount you show off.)}
\end{baihoc}

\begin{ghinhoanh}
\textbf{Short English to remember:}

One real hour beats a week of pretending.

\end{ghinhoanh}

\begin{tuhocnhanh}
1. Em có đang học thật hay đang diễn cảm giác học? \textit{Are you really studying or acting out the feeling of studying?}

2. Làm sao để một giờ học của em thật hơn? \textit{How can your one hour of study become more real?}
\end{tuhocnhanh}
\ngancach

\section{Đừng xin động lực, hãy xin cho mình kỷ luật.}
\begin{truyen}{Đừng xin động lực, hãy xin cho mình kỷ luật.}{Wake up.}
\chuhoa{C}{âu nói nghe cứng nhưng rất thật. Động lực lên xuống như cảm xúc; còn kỷ luật giúp em học được cả những ngày tâm trạng không đẹp.}
\textit{(The line sounds tough but true. Motivation rises and falls like emotion; discipline helps you study even on days when your mood is not good.)}
\end{truyen}

\begin{baihoc}
Thứ cứu em bền nhất không phải cảm hứng, mà là nề nếp.
\textit{(What saves you most steadily is not inspiration, but routine.)}
\end{baihoc}

\begin{ghinhoanh}
\textbf{Short English to remember:}

Seek discipline, not just motivation.

\end{ghinhoanh}

\begin{tuhocnhanh}
1. Em đang phụ thuộc quá nhiều vào động lực không? \textit{Are you depending too much on motivation?}

2. Kỷ luật nào em cần xây trước? \textit{Which discipline do you need to build first?}
\end{tuhocnhanh}
\ngancach

\section{Khi em chịu làm, thầy cô mới giúp được.}
\begin{truyen}{Khi em chịu làm, thầy cô mới giúp được.}{Wake up.}
\chuhoa{T}{hầy cô có thể giải thích, hướng dẫn, động viên, nhưng không ai học hộ em được. Sự giúp đỡ chỉ có đất để mọc khi em chịu tự bước phần đầu tiên.}
\textit{(Teachers can explain, guide, and encourage, but no one can study in your place. Help only has ground to grow on when you are willing to take the first step yourself.)}
\end{truyen}

\begin{baihoc}
Muốn được kéo lên, em phải tự nhấc mình lên trước một chút.
\textit{(If you want to be pulled upward, you must first lift yourself a little.)}
\end{baihoc}

\begin{ghinhoanh}
\textbf{Short English to remember:}

Help works when effort begins.

\end{ghinhoanh}

\begin{tuhocnhanh}
1. Em có đang chờ người khác làm thay phần của mình? \textit{Are you waiting for others to do your part?}

2. Em sẽ chủ động điều gì trước? \textit{What will you take initiative on first?}
\end{tuhocnhanh}
\ngancach

\section{Tự cứu mình trước, rồi người khác mới kéo thêm được.}
\begin{truyen}{Tự cứu mình trước, rồi người khác mới kéo thêm được.}{Wake up.}
\chuhoa{C}{âu này không lạnh lùng, nó rất thực tế. Không ai có thể kéo xa một người cứ nằm yên. Khi em đã chịu xoay người về phía đúng, mọi sự giúp đỡ đều có ý nghĩa hơn.}
\textit{(This sentence is not cold; it is practical. No one can pull far someone who keeps lying still. Once you turn yourself toward the right direction, all help becomes more meaningful.)}
\end{truyen}

\begin{baihoc}
Sự cứu mình đầu tiên là quyết định không tiếp tục buông xuôi.
\textit{(The first self-rescue is the decision to stop giving up.)}
\end{baihoc}

\begin{ghinhoanh}
\textbf{Short English to remember:}

Turn first, then others can pull.

\end{ghinhoanh}

\begin{tuhocnhanh}
1. Em đã quay về hướng đúng chưa? \textit{Have you turned toward the right direction yet?}

2. Quyết định nào của em sẽ là bước tự cứu đầu tiên? \textit{Which decision will be your first act of self-rescue?}
\end{tuhocnhanh}
\ngancach

\section{Không ai cấm em đổi đời, trừ thói quen lười của chính em.}
\begin{truyen}{Không ai cấm em đổi đời, trừ thói quen lười của chính em.}{Wake up.}
\chuhoa{C}{ả lớp nghe xong không cười nhiều nữa. Có những rào cản ngoài đời là thật, nhưng nhiều khi rào cản gần nhất lại nằm ngay trong nếp sống lười và dễ dãi của mình.}
\textit{(After hearing it, the class did not laugh much. Real barriers do exist in life, but often the nearest barrier lies in our own lazy and indulgent habits.)}
\end{truyen}

\begin{baihoc}
Muốn đổi đời, nhiều khi phải đổi nếp sống trước tiên.
\textit{(If you want to change your life, often you must first change your routine.)}
\end{baihoc}

\begin{ghinhoanh}
\textbf{Short English to remember:}

Your habits may be your strongest barrier.

\end{ghinhoanh}

\begin{tuhocnhanh}
1. Thói quen nào đang chặn đường em rõ nhất? \textit{Which habit is blocking your path most clearly?}

2. Nếu bỏ được nó, đời em sẽ khác thế nào? \textit{How would your life change if you removed it?}
\end{tuhocnhanh}
\ngancach

\section{Câu giật mình tốt nhất là câu khiến em đứng dậy ngay.}
\begin{truyen}{Câu giật mình tốt nhất là câu khiến em đứng dậy ngay.}{Wake up.}
\chuhoa{T}{hầy không cần học sinh chỉ gật gù thấy đúng. Điều thầy mong là một câu nói đủ mạnh để sau khi nghe xong, em thật sự bước vào việc phải làm, không chờ thêm một lúc thuận tiện hơn nữa.}
\textit{(The teacher does not need students only to nod and agree. What he hopes for is one sentence strong enough that after hearing it, you actually stand up and do what must be done, without waiting for a more convenient moment.)}
\end{truyen}

\begin{baihoc}
Sự tỉnh ra đẹp nhất là sự tỉnh ra dẫn đến hành động ngay lập tức.
\textit{(The best awakening is the one that leads to immediate action.)}
\end{baihoc}

\begin{ghinhoanh}
\textbf{Short English to remember:}

The best wake-up call gets you moving now.

\end{ghinhoanh}

\begin{tuhocnhanh}
1. Sau câu này, em sẽ đứng dậy làm điều gì đầu tiên? \textit{After this line, what will you stand up and do first?}

2. Em sẽ làm nó ngay hay lại để sau? \textit{Will you do it now or postpone it again?}
\end{tuhocnhanh}
\ngancach